\subsection{実務上の考慮}

\subsubsection{ビジネスケース}

ビジネスケースとは、推奨される業務上の判断を、費用、便益、リスクなど複数の観点からまとめた文書である。投資判断の土台として使われる。単なる説得資料ではなく、意思決定者が判断に必要な情報を整理した成果物である。

\subsubsection{複数通貨分析}

国境をまたぐ意思決定では、為替変動を考慮する必要がある。過去データを使うことも多いが、為替には不確実性がある。したがって、複数通貨の分析では、為替前提、変動幅、リスク対応を明示することが重要である。

\subsubsection{システム思考}

ソフトウェア技術者が関わる環境は複雑である。システム思考は、個別の部品ではなく、全体の関係、相互作用、フィードバック、時間遅れを考える方法である。これにより、ソフトウェアが組織にどのように価値をもたらすかを、より広い視点で説明できる。

\par\smallskip
\begin{center}
\centering
\IfFileExists{assets/generated_figures/ch15/fig-ch15-08.png}{\includegraphics[width=0.96\linewidth]{assets/generated_figures/ch15/fig-ch15-08.png}}{\fbox{\parbox[c][48mm][c]{0.9\linewidth}{\centering 画像生成待ち\\ビジネスケースの構成図}}}
\par\smallskip
\refstepcounter{figure}\label{fig:ch15-08}図 \thefigure: ビジネスケースの構成図
\end{center}
図 \thefigure は、ビジネスケースの構成図を視覚的に整理したものである。
\par\smallskip
